\chapter{Conclusion and Future Work}

\section{Conclusion}

This dissertation implemented a proof-of-concept pipeline for graph-based malware triage from Windows memory forensics: Volatility~3 extraction, typed memory-snapshot graphs, graph-level features, GNN scoring (GIN, GraphSAGE, GAT, GINE), and structured evidence with uncertainty fields \cite{volatility3docs}, \cite{fey2019pyg}.

The primary empirical finding is label--signal mismatch on benign baselines: \textbf{14 of 15} benign-labelled manifest rows are marked uncertain because rule-based triage still reports HIGH or CRITICAL activity. Folder names (\texttt{NoVirus}) are an imperfect proxy for ground truth. This result should be reported openly in any memory-forensics ML study, not treated as a labelling failure to hide.

Secondary contributions are reproducible engineering artefacts: manifests, data contracts, grouped evaluation scripts, tabular baselines, and JSON evidence for analyst review. Together they show that graph-based representation can organise memory-forensics evidence at scale, while detection metrics on $n{=}30$ remain feasibility evidence until larger corpora and archived checkpoints are available.

\section{Answers to Research Questions}

\textbf{RQ1: How can memory artefacts be represented as a graph?} Volatility CSV outputs map to typed nodes and edges (processes, memory regions, DLLs, network objects, intent relations). The manifest records scale and triage metadata per sample.

\textbf{RQ2: Which entities and relationships are most useful?} Process--memory and process--network links are central; intent edges align behaviour with ATT\&CK-style patterns \cite{mitre2026enterprise}. Memory-region nodes dominate graph size in case studies such as WannaCry.

\textbf{RQ3: How can GNNs support detection with readable evidence?} Binary and two-model paths combine learned scores with graph attributes, top nodes, and JSON evidence. Uncertain benign flags motivate analyst-review states rather than forced binary decisions.

\textbf{RQ4: What limitations appear on a small corpus?} High graph-size variance, 14/15 uncertain benign rows, and family pairing require grouped splits and modest claims \cite{sommer2010outside}. Placeholder metric tables document what remains to be run experimentally.

\section{Reflection on the Research Aim}

The aim---to design and evaluate a graph-based deep learning pipeline for memory-forensics malware detection---was met at the pipeline and evidence level. Extraction, graph construction, manifests, models, and analyst-facing outputs are implemented. Full quantitative detection evaluation is partial because checkpoints and \texttt{results\_*.json} files were not archived; Chapter~\ref{ch:evaluation} states that boundary explicitly.

The work does not claim GNNs outperform all alternatives on all malware. It demonstrates feasibility, foregrounds benign-label uncertainty, and specifies the next experiments (grouped CV, tabular comparison, calibration plots).

\section{Contributions}

\begin{itemize}
\item Volatility~3 extraction workflow and artefact contracts;
\item heterogeneous memory-snapshot graph schema;
\item graph-level feature engineering linked to triage;
\item GNN and tabular evaluation scripts with family-grouped CV;
\item structured evidence and uncertainty reporting;
\item dataset manifest for reproducibility;
\item optional dashboard upload-session scanning.
\end{itemize}

\section{Future Work}

\begin{itemize}
\item Expand corpus with clean enterprise baselines and known-good capture times;
\item Run and archive grouped CV for GIN/SAGE/GAT/GINE and tabular baselines; populate Table~\ref{tab:gnn_comparison};
\item Add GNNExplainer \cite{ying2019gnnexplainer} or similar graph explanations;
\item Test on families held out entirely from training;
\item Improve label curation with analyst feedback on uncertain rows;
\item Monitor feature drift across manifest versions;
\item Compare against gradient boosting and operational EDR telemetry \cite{pedregosa2011sklearn}, \cite{ucci2019survey}.
\end{itemize}

\section{Implications for Practice}

Practitioners should treat HIGH or CRITICAL verdicts on benign-labelled baselines as expected when memory contains security tools or noisy host state. Workflow: acquire memory, run plugins, build graphs, inspect triage JSON, then use GNN scores as a secondary signal. Educators can use manifests and figures without requiring enterprise-scale labelled data.

\section{Final Statement}

Graph-based deep learning is a plausible direction for memory-forensics malware triage because system behaviour is relational. The strongest outcome of this project is a reproducible, explainable path from memory dump to evidence---with honest reporting that benign labels and triage signals often disagree. Detection metrics on the current corpus should support feasibility claims only until grouped benchmarks on larger data replace the placeholders in Chapter~\ref{ch:evaluation}.
